% Transcription — fonds Alexandre Grothendieck, shelfmark 156-9, batch 1 (pages 1-20).
% Established from the facsimile archives/batches/156-9/batch-01.pdf (cover sheet
% excluded), per the skill transcribe-grothendieck.
%
% Pass: Opus 5.5 (claude-opus-5-5), 2026-10-03 — first pass, unchecked against the pages by a human.
%
% All twenty pages are transcribed; none is skipped.
%
% Page 1 is headed « GF IX » (géométrie des formes, chapter IX) and carries a
% graph of the types of « ateliers » with a fast prose paragraph; page 2
% continues it (dated 15.7 in the margin). Page 3, dated 15 July 86, is the
% clean « Diagramme des principaux types d'ateliers » with its legend.
% Pages 4-20 are his own run paginated 1-17, dated 5 July: the magasin
% M = M_0 ⨿ M_1 attached to an ordered set L (vertices and intervals I_{a,b}),
% its order ≤, the relations ◁, ⊴, ≪°, |o|, the result that the prefigures are
% the finite chains, the infinite intervals I_{<a}, I_{>a} and the widened
% magasin M̂, the complementary prefigure F*, the theorem
% cosupp°(F) = Omb°(F*), then the « autre présentation » on a combinatorial
% oriented segment (I, ∂I) with the classification of vertices (isolated,
% lacunary, evanescent, semi-evanescent) and their exchange under duality.
% The run continues past page 20.
%
% The hand is fast throughout; the formulas carry and much connecting prose
% does not. The negation of a relation is a bar over its sign, rendered with
% \overline. ◁ is \lhd, ⊴ \trianglelefteq, the relation he writes as « ≪ with
% a ring over it » \mathring{\ll}.
%
% The dating is the inventory's own for the shelfmark.
\documentclass[11pt,a4paper]{article}
\input{../preamble/grothendieck}

\folder{156-9}
\batch{1}
\pages{1}{20}
\foldertitle{[Chapitre] IX et IX bis. [Ateliers] : notes manuscrites (05-15/07/1986).}
\dating{1986}
\watermark{Édition de démonstration}

\begin{document}

\section{GF IX}

\page{1}
\note{En tête de la feuille, de sa main : « GF IX » (IX souligné et surligné), puis au crayon : « (voir p.~75 la mise au pt sur la terminologie) ».}

\drawing{un graphe des types d'ateliers, reliés par des traits et des flèches. En haut, « Atelier général $(\mathcal{A};\trianglelefteq,\mathring{\ll},|o|)$ » cerclé ($\mathcal{A}$) ; sous lui « Atelier divisible » (souligné, traversé par le trait qui monte vers l'atelier général). Au niveau suivant : « atelier spatial $(\mathcal{A};\trianglelefteq,|o|)$ » avec $\Sigma$ cerclé, relié à l'atelier général ; « atelier local $(\mathcal{A};\trianglelefteq,\mathring{\ll},|o|_{\mathcal{L}})$ » ; « atelier fidèle $(\mathcal{A};\trianglelefteq,\mathring{\ll})$ », relié à l'atelier général. Plus bas, « atelier ponctuel $\Sigma\simeq\mathfrak{P}(\mathcal{L})$ », relié à l'atelier local et à l'atelier fidèle, et par un trait tireté à l'atelier spatial ; une flèche de l'atelier spatial et une flèche de l'atelier ponctuel descendent vers ($\mathcal{L}$) « atelier ensembliste », d'où une flèche vers « at. ens. modéré », puis vers « at. discret », puis vers « at.-point » (au-dessus, un mot biffé, \ill{}). À droite, « atelier modéré » (un premier « atelier modéré » est biffé plus haut) reçoit une longue courbe venue de l'atelier général et une autre de l'atelier fidèle, avec au crayon « at. mod ? » ; une flèche tiretée en part vers « at. ens. modéré » ; « maquette $\Sigma\simeq\mathfrak{P}(\mathcal{A})$ », encadrée au crayon, est reliée par une double flèche à l'atelier modéré et envoie une flèche vers « at. discret ». Un trait horizontal orangé sépare « at. discret » et « at.-point » du reste.}

\marginal{NB Si les ateliers mod. sont fidèles, le diagramme se simplifie : at. mod. se place entre at. fid. et maquette.}

Je m'\uncertain{intéresserai} plus particulièrement \add{\uncertain{\ill{}}} par les ateliers \emph{spatiaux}* (\uncertain{ça va de soi !}), incluant les ateliers ensemblistes, comme ceux définis par des \struck{ens. ordonnés} « intervalles \uncertain{ordonnés} », et les maquettes (at. spatiaux), \ill{} \uncertain{qui ont} par certains côtés, \ill{}. Parmi ceux-ci, pour un travail \uncertain{modéré} (\uncertain{dans} l'esprit de la topologie \ill{} ensembliste), je suis \uncertain{intéressé} par les ateliers qui sont de plus locaux. Cela \uncertain{incluent} les maquettes, parmi les exemples précédents, \uncertain{sauf} les maquettes discrètes).

En \uncertain{direction} d'une topologie modérée, je m'intéresserai aux \emph{ateliers modérés} (\uncertain{ensemblistes} ou non) --- qui incluent les maquettes. Je présume qu'il y a lieu de \struck{\uncertain{imposer}} pour les ateliers modérés \ill{} \uncertain{spatialité} (\uncertain{c'est} \ill{} par \uncertain{nécessité} locaux), et \ill{}

\note{En bas à droite : « TSVP ». La lecture de ce paragraphe, écrit vite et chargé de reprises interlinéaires, est très incertaine.}

\page{2}
\note{La page commence au milieu d'une phrase, qui continue le recto (« TSVP », p.~1).}

aussi que l'ensemble $\Sigma$ \add{ordonné} des supports soit « modulaire », i.e.
\[
  A \wedge \operatorname{Sup}_i B_i = \operatorname{Sup}_i A \wedge B_i ,
\]
dans une « algèbre de Boole ». Il est possible que cela implique que l'atelier soit aussi fidèle.

\marginal{À vérifier ? \uncertain{Non},}

\emph{Question.} Si $x \in \mathcal{L}$, alors $x^{\circ} = \operatorname{Supp}^{\circ} x = \operatorname{Supp} x \in \Sigma$ est-il un objet minimal de $\Sigma^{*} = \Sigma \smallsetminus \{0_{\Sigma}\}$ ? En d'autres termes, si $X \in \mathcal{A}$, \struck{$X \in \mathcal{L}$}, et $\operatorname{Supp}^{\circ} X \subset \operatorname{Supp}^{\circ} x$, a-t-on $X = x$ ? En d'autres termes,
\[
  X \neq x \overset{?}{\Longrightarrow} \operatorname{Supp}^{\circ}(X) \not\subset \operatorname{Supp}^{\circ} x ,
  \quad \text{i.e.} \quad \operatorname{cosupp}^{\circ} x \not\subset \operatorname{cosupp}^{\circ} X ,
\]
i.e. $\exists\, Z \in \mathcal{A}$ avec $Z \mathrel{|o|} x$, $Z \mathrel{\overline{|o|}} X$ ? Si cela \ill{} --- revient au cas où $X = y \in \mathcal{L}$ (locaux \ill{}), donc à voir que si $x, y \in \mathcal{L}$, $y \neq x$,
\[
  \exists\, Z \in \mathcal{A} \quad \text{avec} \quad Z \mathrel{|o|} x ,\ Z \mathrel{\overline{|o|}} y .
\]
(Car $x^{\circ} = \operatorname{Sup}_{y} y^{\circ}$, $y \in \mathcal{L}$ t.q.\ $y \subset X$, i.e.\ $y^{\circ} \subset X^{\circ}$.) Je \ill{} pas \uncertain{que} \struck{dans} le cas $x \mathrel{\overline{|o|}} y$.
\note{Dans ce dernier paragraphe la négation de $|o|$ est notée par une barre au-dessus du signe ; le sous-indice de $\operatorname{Sup}$ est écrit sur trois lignes et sa lecture est incertaine.}

\marginal{15.7.}
\emph{Remarques :} Opération de restriction d'un atelier à un contexte, faisable dans les ateliers, \ill{} \uncertain{ou dans} \uncertain{celui} des ateliers spatiaux sans plus ? J'en doute !

\drawing{un graphe au crayon et à l'encre, de sommets S (en haut), LS, M (cerclés), SF, E, MF, EM (cerclé) et Mq (en bas) : S est relié à M par un arc épais et à LS ; SF est au centre, relié à S, à E et à MF ; M est relié à MF par un arc ; MF est relié à EM et à Mq ; LS est relié à E puis à EM. À côté : « DEQFLMPS » et dessous « Mq, Div, v, p ».}

\page{3}
\section*{Diagramme des principaux types d'ateliers (ou de « contextes »)}

15 juillet 86

\drawing{un graphe orienté, les flèches montant vers le haut (traits doublés à l'encre brune), avec hachures au crayon dans deux bandes. Au sommet G $(\mathcal{A})$. Au-dessous, à gauche S $(\Sigma$ ou $\mathcal{A})$ cerclé, au milieu F, à droite L ; sur l'arête de L à G un point « Div » ; sur l'arête de SF à S une marque « [S Div] ». Au niveau suivant SF, LS $(\Sigma$ ou $\mathcal{L})$ (cerclé) et FL $=$ P. Puis FLS $=$ E $(\mathcal{L})$ au centre et M (cerclé) à gauche. Puis Mq (sous M) et EM (cerclé, sous E). En bas D, recevant deux flèches venues de v et de p. Flèches : S, F, L $\to$ G ; SF $\to$ S et F ; LS $\to$ S et L ; P $\to$ F et L ; E $\to$ SF, LS, P ; M $\to$ SF ; Mq $\to$ M ; EM $\to$ M et E ; D $\to$ Mq et EM ; v, p $\to$ D. Pointillés de S à Div et de F à L.}

\emph{Légende :}
\begin{itemize}
  \item G $=$ ateliers généraux
  \item S $=$ --- spatiaux
  \item F $=$ --- fidèles
  \item L $=$ --- locaux
  \item SF $=$ --- spatiaux \uncertain{locaux}\note{sic, semble-t-il ; on attendrait « fidèles ».}
  \item FL $=$ P $=$ --- fidèles locaux \add{ou \uncertain{ponctuaires}}
  \item LS $=$ --- locaux spatiaux (NB LS $=$ S $\cap$ Div, at. spatiaux divisibles)
  \item FLS $=$ E $=$ --- fidèles locaux spatiaux \add{ou ensemblistes}
  \item M $=$ --- modérés
  \item EM $=$ --- ensemblistes (modérés)
  \item Mq $=$ --- maquettes
  \item D $=$ --- discrets
  \item Div $=$ --- (infiniment) divisibles ; SDiv $=$ --- spatiaux (inf.) divisibles
  \item p $=$ \struck{ateliers} ponctuels ($\operatorname{card} \mathcal{A} = 1$)
  \item v $=$ --- vide ($\operatorname{card} \mathcal{A} = 0$).
\end{itemize}
\note{Une flèche marquée « ? » relie LS à FLS $=$ E dans la légende.}

Les huit « carrés élémentaires »

\begin{tikzcd}[column sep=small, row sep=small]
 & d & \\
b \arrow[ur] & & c \arrow[ul] \\
 & a \arrow[ul] \arrow[ur] &
\end{tikzcd}

figurant dans ce diagramme sont des « carrés-intersection », i.e.\ $a = b \wedge c$. On en conclut p.~ex., par « chasse », que $\mathrm{LS} \cap \mathrm{M} = \mathrm{EM}$ ($= \mathrm{E} \cap \mathrm{M}$).

\struck{\ill{}}

Les cadres privilégiés par excellence pour une « géométrie des formes » dans un « contexte » \add{($\doteq \mathrm{S} \cap \mathrm{Div}$)} me semblent à présent : LS (ateliers locaux \struck{\uncertain{ou mieux, spatiaux divisibles}} spatiaux), \struck{\ill{}} qui semblent adéquats pour l'étude des \struck{\ill{}} « \add{cordes et} formes » (pseudo-simplexes topologiques), S (ateliers spatiaux) adéquats (en général) pour les contextes finis), M (ateliers modérés, pour la topologie modérée) et EM (topologie modérée ensembliste).

\marginal{À vérifier que la bonne notion d'\emph{atelier modéré} doit bien impliquer qu'il soit, non seulement \emph{spatial}, mais encore \emph{fidèle}. Sinon, diagramme moins joli ! (Cf.\ \uncertain{ci-contre}, feuille \uncertain{suivante}\ldots)}

\marginal{Je note que l'accent principal s'est nettement déplacé depuis les approches originelles via L, vers des approches subordonnées à S (les \uncertain{originelles} \uncertain{venues} dans la trilogie FLS).}

\page{4}
\note{Pagination de l'auteur : 1 (un premier numéro biffé).}

\emph{5 juillet.} L'affaire a tendance à s'enliser, et j'ai envie de tout reprendre, d'une façon \struck{critique} un peu différente.

NB. $\mathcal{L}$ est un ens. ordonné ; je lui associe un magasin \struck{\ill{}}
\[
  \mathfrak{M} = \mathfrak{M}_0 \amalg \mathfrak{M}_1 ,
\]
$\mathfrak{M}_0 = \mathcal{L}$, ensemble des « \emph{sommets} » virtuels, ou « \emph{lieux} » ; $\mathfrak{M}_1 \overset{\simeq}{=} \operatorname{Drap}_2(\mathcal{L})$. Si $\{a,b\}$ ($a<b$) est dans $\operatorname{Drap}_2(\mathcal{L})$, je note $I_{\{a,b\}}$ ou $I_{a,b}$ ou $I_{b,a}$ la structure \uncertain{associée}, appelée aussi \emph{intervalle} déterminé par $\{a,b\}$.

\marginal{\uncertain{canonique} \ill{} \uncertain{associé} à $\mathcal{L}$}

\struck{Relation d'incidence : $X \lhd Y$ ssi $X \in$}

Notation $\partial X$, $X \in \mathfrak{M}$ :
\[
  \text{(a)} \quad
  \begin{cases}
    \partial I_{\{a,b\}} = \{a,b\} \subset \mathfrak{M}_0 \\
    \partial x = \emptyset \subset \mathfrak{M}
  \end{cases}
  \qquad
  \text{(b)} \quad
  \begin{cases}
    \delta I_{a,b} = \partial I_{a,b} = \{a,b\} \\
    \delta x = \{x\}
  \end{cases}
\]

\emph{Relation d'incidence :} $X \lhd Y$ ssi $X = x \in \mathcal{L}$, $Y \in \mathfrak{M}_1$, et $x \in \partial Y$. Donc la condition $\partial X$ de la \uncertain{relation} est \uncertain{celle} d'incidence :
\[
  \partial X = \{ Y \in \mathfrak{M} \mid Y \lhd X \} \qquad (X \in \mathfrak{M}) .
\]

[Il semble qu'il vaut mieux introduire d'abord la relation $\leq$ sur $\mathfrak{M}$, \uncertain{d'interpréter} \struck{\ill{}} $I_{ab}$, $I_{<a}$, $I_{>b}$ \uncertain{comme} des « intervalles » dans $\mathfrak{M}$, \ill{} \uncertain{via} \ill{}.]

\marginal{Écrit en oblique dans la marge gauche, de lecture très incertaine : \ill{} $[I_{a,b}] = \{x \in \mathcal{L} \mid a \leq x \leq b\}$, \ill{} $]a,b[$ \ill{} $x = $ \ill{} ; $\operatorname{c} \| I_{a,b} | a | = $ \ill{} $X \in \mathfrak{M}_1$ \ill{} ; \ill{} $\delta X \lhd \partial X \subset \delta X$.}

\emph{Relation $\ll$.} On a par définition
\[
  X \mathrel{\mathring{\ll}} Y \iff |X| \subset |Y| \ \text{et} \ X \mathrel{\overline{\lhd}} Y .
\]
\[
  X \mathrel{\mathring{\ll}} Y \iff
  \begin{cases}
    X = Y \\
    \text{ou } X = x \in \mathcal{L},\ Y = I_{a,b},\ a < x < b \\
    \text{ou } X = I_{a,b},\ Y = I_{a',b'},\ a' \leq a < b \leq b'
  \end{cases}
\]
\note{La barre au-dessus de $\lhd$ note la négation. Un mot biffé et noirci précède « $X \overline{\lhd} Y$ ».}

\marginal{Relation d'ordre. Mag 2, OK \note{encadré, avec une flèche qui descend vers le bas de la page}}

\emph{Remarque.} $X \mathrel{\mathring{\ll}} Y \Longrightarrow \delta X \cup \delta Y \in \operatorname{Drap}(\mathcal{L})$.

\emph{Conséquence.} Relation $X \ll Y$ $\bigl(\overset{\mathrm{def}}{\iff} \exists\, Y',\ X \mathrel{\mathring{\ll}} Y' \trianglelefteq Y\bigr)$ :
\[
  X \ll Y \iff X \lhd Y \ \text{ou} \ X \mathrel{\mathring{\ll}} Y \iff |X| \subset |Y|
\]
\note{Sous « $X \lhd Y$ ou $X \mathring{\ll} Y$ », une accolade : « \struck{contradictoires} mutuellement exclusifs ». Au début de la ligne, un mot biffé, \ill{}.}
\[
  \iff
  \begin{cases}
    X = Y \\
    \text{ou } X = x \in \mathcal{L},\ Y = I_{a,b},\ a \leq x \leq b \\
    \text{ou } X = I_{a,b},\ Y = I_{a',b'},\ a' \leq a < b \leq b'
  \end{cases}
\]

\page{5}
\note{Pagination de l'auteur : 2.}

Avant d'introduire $|o|$, on introduit relation d'ordre sur $\mathfrak{M}$ :
\[
  X < Y \overset{\mathrm{def}}{\iff}
  \begin{cases}
    X = x,\ Y = y,\ x, y \in \mathcal{L},\ x < y \\
    X = x,\ Y = I_{a,b},\ x \leq a < b \\
    X = I_{a,b},\ Y = y,\ a < b \leq y \\
    X = I_{a,b},\ Y = I_{a',b'},\ a < b \leq a' < b'
  \end{cases}
\]
\[
  X < Y \iff \bigl( X \neq Y \ \text{et} \ (\forall\, x \in \delta X,\ y \in \delta Y,\ \text{on a } x \leq y) \bigr)
  \iff \bigl( |X| \leq |Y| \bigr)
\]
\note{Dans cette seconde ligne le signe $<$ est repassé d'un trait épais ; au-dessus de « $|X| \leq |Y|$ » est ajouté « $|X| \neq |Y|$ et », relié par un trait à « $X \neq Y$ et ».}

\marginal{NB \ill{} pour deux parties $A$, $B$ \ill{} (ou relations) d'un ens. ordonné, $A < B$ si \ill{} ($A \leq B \iff A = B$ ou $A < B$) i.e.\ $A \leq B$, $A < B$ \ill{} \note{écrit en oblique, lecture très incertaine}}

\emph{Remarques.} $X < Y \Longrightarrow \delta X \cup \delta Y \in \operatorname{Drap}^{*}(\mathcal{L})$.

Notation pour $D \in \operatorname{Drap}^{*}(\mathcal{L})$ :
\begin{itemize}
  \item $\operatorname{or} D = $ plus petit élément de $D$,
  \item $\operatorname{ex} D = $ --- grand ---
\end{itemize}
et \struck{$\operatorname{or} X$} pour $X \in \mathfrak{M}$ :
\[
  \operatorname{or} X = \operatorname{or} \delta X , \qquad \operatorname{ex} X = \operatorname{ex} \delta X .
\]

\emph{Remarques.} a) $X < Y \Longrightarrow (\operatorname{ex} X \leq \operatorname{or} Y)$ \struck{et $\operatorname{or} X \leq \operatorname{ex} Y$}

b) $X \leq Y \Longrightarrow (\operatorname{or} X < \operatorname{ex} Y$, \emph{ou} $X = Y = a \in \mathcal{L})$

\struck{c) $X < Y$, $X' \mathrel{\mathring{\ll}} X$, $Y' \mathrel{\mathring{\ll}} Y \Longrightarrow X' < Y'$}

\struck{On a : si $X < Y$, $X' \ll X$, $Y' \ll Y$, $X', Y' \in \mathfrak{M}$,}
\struck{$X' \leq Y'$, sauf si $X' = \operatorname{ex} X = Y' = \operatorname{or} Y$}

\struck{Définition : $X \mathrel{|o|} Y \overset{\mathrm{def}}{\iff}$}
\struck{$X < Y$ ou $Y < X$, i.e.\ $\{X,Y\} \in \operatorname{Drap}(\mathfrak{M})$}
\note{Ces trois lignes sont barrées de traits obliques.}

\emph{Proposition.} Soient $X, Y \in \mathfrak{M}$, $X', Y' \in \mathfrak{M}$, $X' \ll X$, $Y' \ll Y$. \emph{Alors} \struck{\ill{}}
\[
  X < Y \Longrightarrow X' \leq Y' , \ \struck{et}
\]
et on a même
\[
  X < Y \Longrightarrow \bigl( X' < Y' \ \text{ou} \ X' = \operatorname{ex} X = Y' = \operatorname{or} Y \bigr) .
\]

Dans ce dernier cas, on aura $X' \trianglelefteq X$, $Y' \trianglelefteq Y$, donc on ne peut avoir $X' \mathrel{\mathring{\ll}} X$, $Y' \mathrel{\mathring{\ll}} Y$ que si $X' = X$, $Y' = Y$, \add{ce qui} \struck{\uncertain{contredit l'hyp.}} implique $X = Y$, et est \struck{exclu} \uncertain{contraire} à : $X < Y$. Donc

\emph{Cor.} $X < Y$, $X' \mathrel{\mathring{\ll}} X$, $Y' \mathrel{\mathring{\ll}} Y \Longrightarrow X' < Y'$.

\emph{Déf.}
\[
  X \mathrel{|o|} Y \overset{\mathrm{def}}{\iff} X < Y \ \text{ou} \ Y < X \quad \text{i.e.} \quad \{X, Y\} \in \operatorname{Drap}_2(\mathfrak{M}, \leq) .
\]

Le corollaire précédent montre Mag 3, et Mag 4 immédiat, car $a < I_{ab} < b$, donc $X \lhd Y \Longrightarrow X \mathrel{|o|} Y$.

\page{6}
\note{Pagination de l'auteur : 3 (un premier numéro biffé).}

\struck{$X \asymp Y \iff$}

Soient \ill{} \uncertain{l'équivalence} \uncertain{générale} \uncertain{en} $\asymp$,

\struck{On \uncertain{voit} que pour $X \neq Y$, on a $X$}
\struck{On veut que pour $X \neq Y$, on \uncertain{puisse avoir}}
\struck{$\widetilde{X} \cap \widetilde{Y} \neq \emptyset$ que si \ill{} $\delta$}

\[
  X \neq Y \Longrightarrow \widetilde{X} \cap \widetilde{Y} = \delta X \cap \delta Y , \ \text{de cardinal } 0 \text{ ou } 1
\]

\struck{et si $\widetilde{X} \cap \widetilde{Y} \neq \emptyset$, on est dans la situation $X, Y \in \mathfrak{M}_1$, $\operatorname{card}(\delta X \cap \delta Y) = 1$,}
\struck{$X \lhd Y$ ou $Y \lhd X$ ou}

\struck{Donc $X \asymp Y \iff X = Y$ ou $X \lhd Y$ ou $Y \lhd X$}
\note{Tout le début de la page, jusqu'ici, est rayé de traits horizontaux et encadré d'un trait à gauche. Le signe noté ici $\asymp$ est tracé comme un $\asymp$ barré ; on le rend partout par $\asymp$.}

\emph{Proposition.}
\[
  X \asymp Y \iff X \overset{\sim}{=} Y \ \text{ou} \ X \mathrel{|o|} Y \quad \text{i.e.} \quad \{X, Y\} \in \operatorname{Drap}(\mathfrak{M})
\]
\note{Le signe d'égalité porte un petit arc au-dessus ; lecture incertaine.}

\emph{\struck{\ill{}} Dém.} OPS $X \neq Y$. Alors
\[
  \widetilde{X} \cap \widetilde{Y} = \delta X \cap \delta Y \ \text{est de cardinal } 0 \text{ ou } 1 .
\]
\uncertain{On distingue} les deux cas.

a) $\widetilde{X} \cap \widetilde{Y} = \emptyset$. Alors la définition \add{de $X \asymp Y$} signifie $\widetilde{X} \mathrel{|o|} \widetilde{Y}$, et on voit que cela équivaut \uncertain{aussi} : $X \mathrel{|o|} Y$, par l'hypothèse \struck{\ill{}} $\widetilde{X} \cap \widetilde{Y} = \emptyset$.

b) $\operatorname{card} \widetilde{X} \cap \widetilde{Y} = 1$, i.e.\ $\widetilde{X} \cap \widetilde{Y} = \{b\}$. \struck{Alors Donc} \struck{$X \asymp Y$ signifie} Si $X \lhd Y$ ou $Y \lhd X$, alors $X \asymp Y$ \uncertain{et} $X \mathrel{|o|} Y$, \uncertain{et} la prop.\ est satisfaite. OPS \struck{que $X$ et $Y$} \ill{} \add{$X, Y \notin \widetilde{X} \cap \widetilde{Y}$}, \struck{\ill{}} mais alors $X, Y \in \mathfrak{M}_1$, \struck{\ill{}} et $X \asymp Y$ \struck{signifie} \uncertain{implique} $X \mathrel{|o|} Y$, \struck{\ill{}} \uncertain{donc} (\uncertain{par construction}) $X = I_{ab}$, $Y = I_{b,c}$, $a < b < c$. Mais alors $(\widetilde{X} \smallsetminus \{b\}) \mathrel{|o|} (\widetilde{Y} \smallsetminus \{b\})$ i.e.\ $X \asymp Y$, cqfd.

\emph{Corollaire.} Pour qu'une partie \add{finie} $\mathfrak{F}$ de $\mathfrak{M}$ soit une préfigure de $\mathfrak{M}$ (i.e.\ $\mathfrak{F}$ soit une figure) il faut et il suffit qu'elle soit totalement ordonnée par $\leq$.

\page{7}
\note{Pagination de l'auteur : 4.}

\emph{NB} « Figures » signifie « figures finies » --- i.e.\ \add{$\mathfrak{F} \subset \mathfrak{P}_f(\mathfrak{M})$} on travaille dans l'atelier des figures finies. On trouve ainsi
\[
  \operatorname{Préfig}(\mathfrak{M}) = \operatorname{Drap}^{*}(\mathfrak{M}, \leq)
\]
\[
  \operatorname{Fig}(\mathfrak{M}) = \operatorname{Drap}^{*}(\mathfrak{M}, \leq) \cap \mathfrak{P}_f(\mathfrak{M}, \trianglelefteq)
\]
Les figures sont les parties finies de $\mathfrak{M}$, totalement ordonnées par $\leq$, et fermées pour $\trianglelefteq$.

Pour l'étude des supports contractibles et de l'\uncertain{utilisation} du complémentaire, il sera utile d'introduire, pour toute préfigure, une « préfigure complémentaire ». Comme $\mathcal{L}$ n'a pas de plus petit \emph{et} de plus grand élément, ceci fait sortir des préfigures proprement dites, en nous faisant introduire des « intervalles infinis » $I_{<a}$, $I_{>a}$. Même si $\mathcal{L}$ a plus petit et plus grand élément, l'introduction de $I_{<a}$, $I_{>b}$ rend la formation des préfigures \uncertain{complémentaires} plus élégante ; on obtient la complémentaire $\mathfrak{F}'$ d'une préfigure $\mathfrak{F}$ comme préfigure, \struck{\ill{}} avec $\delta \mathfrak{F}' \subset \delta \mathfrak{F}$, alors que \struck{\uncertain{généralement}} dans le contexte des préfigures « ordinaires », la définition \uncertain{imposerait} \uncertain{des} \ill{} : on aurait seulement
\[
  \delta \mathfrak{F}' \subset \delta \mathfrak{F} \cup \{\omega_0, \omega_1\}
\]
(\uncertain{avec} plus petit et plus grand élément). (De façon précise, on aurait $\mathcal{D} = (\mathfrak{F} \cap \mathcal{D}) \amalg (\mathfrak{F}' \cap \mathcal{D})$ --- tout point de $\mathcal{D}$ (extr.\ de $\mathcal{L}$) --- \uncertain{appartient} à une des deux préfigures
\note{La lettre lue ici $\mathcal{D}$ est incertaine.}

\page{8}
\note{Pagination de l'auteur : 5.}

$\mathfrak{F}$, $\mathfrak{F}'$, et une seule). \struck{Ainsi} Quand on se limite aux préfigures \add{(\uncertain{généralisées})} \uncertain{irrédondantes}, \uncertain{dans} le \uncertain{cadre} de l'introduction des $I_{<a}$, $I_{>a}$, \ill{} $\delta\mathfrak{F} = \delta\mathfrak{F}'$, alors que \ill{} \ill{} pour \ill{} la formalisme \ill{} \uncertain{généralisée}, on \uncertain{a} les préfigures \add{de $<$} \ill{} les extrémités \uncertain{\ill{}} \ill{} communes.

Introduction de nouveaux éléments, pour former $\widehat{\mathfrak{M}} = \widehat{\mathfrak{M}}_0 \amalg \widehat{\mathfrak{M}}_1$ :
\[
  \widehat{\mathfrak{M}}_0 = \mathfrak{M}_0 = \mathcal{L} , \qquad
  \widehat{\mathfrak{M}}_1 = \mathfrak{M}_1 \amalg \mathfrak{M}_1^{-} \amalg \mathfrak{M}_1^{+}
\]
\[
  \mathfrak{M}_1^{-} = \{ I_{<a} \mid a \in \mathcal{L},\ a \text{ non minimal dans } \mathcal{L} \}
\]
\[
  \mathfrak{M}_1^{+} = \{ I_{>a} \mid a \in \mathcal{L},\ a \text{ non maximal dans } \mathcal{L} \}
\]

$I_{<a}$, $I_{>a}$ \struck{pour} sont ici des symboles graphiques, auxquels pour l'instant on n'associe aucune interprétation. Plus \uncertain{intuitivement}, \struck{\ill{}} $I_{<a}$ figure \uncertain{l'ensemble} \ill{} formé des $x \in \mathcal{L}$, $x < a$, ou \uncertain{mieux}, les $X \in \mathfrak{M}$ tels que $X < a$, et symétriquement pour $I_{>a}$.

On pose
\[
  \partial I_{<a} = \delta I_{<a} = \text{\struck{$\ldots$}} = \{a\}
\]
\[
  |I_{<a}| \overset{\mathrm{def}}{=} \mathcal{L}_{\leq a} = \{ x \in \mathcal{L} \mid x \leq a \}
\]
\[
  |I_{<a}|^{\circ} \overset{\mathrm{def}}{=} |I_{<a}| \smallsetminus \partial I_{<a} = \mathcal{L}_{<a} = \{ x \in \mathcal{L} \mid x < a \}
\]
\note{Dans la première ligne, deux formules intermédiaires sont biffées par des hachures ; dans la troisième, un mot biffé précède $|I_{<a}|$.}
et de même pour $I_{>a}$. Cela étend $\trianglelefteq$ et $\lhd$ à $\widehat{\mathfrak{M}}$. On \uncertain{étend} pareillement $\mathring{\ll}$, en \uncertain{écrivant} les nouveaux cas de $\mathring{\ll}$, \uncertain{indiquant} $I_{<a}$

\page{9}
\note{Pagination de l'auteur : un chiffre lu 8, peut-être 6 ; la page continue la précédente.}

\[
  \begin{cases}
    I_{<a} \mathrel{\mathring{\ll}} I_{<b} \iff a \leq b \\
    X \mathrel{\mathring{\ll}} I_{<a} \ \text{si} \iff X \leq a \quad \text{(cf.\ déf.\ldots)}
  \end{cases}
\]
\note{Après « $X \leq a$ » quelques signes sont biffés par des hachures, et « dans \ill{} » est biffé sous « (cf.\ déf.\ldots) » ; le $\leq$ est repassé.}
(où $X \in \mathfrak{M}$), et \uncertain{analogue} pour \struck{$I_{>}$} les relations impliquant les intervalles infinis supérieurement.

\marginal{\ill{} $X \mathrel{\mathring{\ll}} Y \iff |X| \subset |Y|$ \uncertain{et} $X \mathrel{\overline{\lhd}} Y$ \ill{} pour $X, Y \in \widehat{\mathfrak{M}}$ \ill{} \note{écrit en oblique dans la marge gauche ; lecture très incertaine}}

On notera :

a) On n'a jamais une relation
\[
  I_{<a} \mathrel{\mathring{\ll}} X \qquad \text{pour } X \in \mathfrak{M} .
\]
Par exemple, si $\mathcal{L}$ a un plus petit él.\ $\omega_0$, on n'a \emph{pas} $I_{<a} \mathrel{\mathring{\ll}} I_{\omega_0,a}$, \struck{car (\ill{} intuitivement)} \ill{} \uncertain{deux fois}. $I_{\omega_0,a} \mathrel{\mathring{\ll}} I_{<a}$. En fait, $I_{\omega_0,a}$ doit être \ill{} \uncertain{définissant} une \uncertain{modalité}\ldots) comme une \emph{\uncertain{substitution}} \uncertain{effective} de $I_{<a}$. (Si on \struck{avait} \add{\uncertain{admettait}} les deux relations : \ill{}, on \uncertain{aurait} $I_{<a} = I_{\omega_0,a}$, ce qui n'est pas \uncertain{conforme} aux intentions présentes \ill{} \uncertain{ni} \struck{\ill{}} $I_{<a} \ll I_{\omega_0,a}$, mais \uncertain{la relation inverse} \struck{\ill{}} $I_{\omega_0,a} \ll I_{<a}$, et $|I_{\omega_0,a}| = |I_{<a}|$.)

b) On n'a jamais une relation $I_{<a} \mathrel{\mathring{\ll}} I_{>b}$, ni l'inverse.

\marginal{Mag 2 OK \uncertain{sur} $\widehat{\mathfrak{M}}$}

Relation d'ordre $\leq$ sur $\widehat{\mathfrak{M}}$. On pose, \uncertain{comme}, \struck{si $X, Y \in \widehat{\mathfrak{M}}$} \uncertain{étend} $\leq$ à $\widehat{\mathfrak{M}}$ par :

\struck{$X \leq Y \iff \forall\, x \in \delta X,\ \forall\, y \in \delta Y,\ x \leq y$}

(on écrit aussi $\delta X \leq \delta Y$, relation qui a un sens pour deux parties $A$, $B$ quelconques de $\mathcal{L}$ : $A \leq B$)
\note{La ligne de définition est biffée d'un trait horizontal et la parenthèse qui suit de traits obliques.}

\page{10}
\note{Pagination de l'auteur : 7.}

\struck{Relation $|o|$ sur $\widehat{\mathfrak{M}}$} \add{($X \neq Y$ \uncertain{et})}
\[
  X < Y \iff |X| \leq |Y| \qquad \text{On trouve}
\]
\[
  \text{\struck{$X \mathrel{|o|} Y \iff X < Y$ ou $Y < X$, i.e.\ $\{X,Y\} \in \operatorname{Drap}_2(\widehat{\mathfrak{M}}, \leq)$}}
\]
\note{La première ligne, avec l'ajout encerclé « $X \neq Y$ et » muni d'un signe $\exists$ au-dessus, n'est pas toute biffée ; lecture incertaine de cette articulation.}

\[
  X \leq Y \iff
  \begin{cases}
    1^{\circ})\ \delta X \leq \delta Y \\
    2^{\circ})\ \text{si $X$ ou $Y$ est de la forme $I_{<a}$, $I_{>b}$, on exige \struck{pour} ceci :} \\
    \qquad \text{a) si } Y = I_{<a},\ \text{alors } X = Y \\
    \qquad \text{b) si } X = I_{>a},\ \text{alors } Y = X
  \end{cases}
\]
\note{Dans a), « donc $X = Y$ » est biffé avant « alors $X = Y$ » ; sous les deux cas une ligne de signes est biffée par des hachures. Le $\leq$ et le $<$ de la première ligne de la page sont repassés.}

\emph{NB} On a dans tous les cas
\[
  X \leq Y \Longrightarrow |X| \leq |Y| , \qquad X < Y \Longrightarrow |X| < |Y| .
\]

\emph{NB} \struck{\ill{}} On n'a jamais $I_{>a} \leq I_{<b}$ à cause de $2^{\circ}$ a) (ou b)) ci-dessus. On a, \uncertain{bien sûr},
\[
  I_{<a} \leq I_{>b} \iff a \leq b .
\]
\note{Dans la formule, le $>$ de $I_{>b}$ est repassé sur un autre signe ; un trait ponctué double le $\leq$.}

On pose \struck{\ill{}}, pour $X, Y \in \widehat{\mathfrak{M}}$
\[
  X \mathrel{|o|} Y \iff X < Y \ \text{ou} \ Y < X \quad \text{i.e.} \quad \{X, Y\} \in \operatorname{Drap}_2(\widehat{\mathfrak{M}}, \leq) .
\]

Vérification des axiomes. Mag 1 OK, Mag 2 OK, Mag 3 OK car
\[
  X < Y,\ X' \mathrel{\mathring{\ll}} X,\ Y' \mathrel{\mathring{\ll}} Y \Longrightarrow X' < Y' .
\]
\[
  \struck{\ill{}} \quad X < Y,\ X' \ll X,\ Y' \ll Y \Longrightarrow X' \leq Y'
\]
\[
  \hookrightarrow X' < Y' \ \text{ou} \ X' = \operatorname{ex} X = Y' = \operatorname{or} Y .
\]
\note{Une accolade sous les hypothèses de la deuxième ligne renvoie, par une flèche, à la troisième ; au-dessus de celle-ci, « $\leftarrow X' =$ » est biffé.}

\emph{Mag 4} clair.

\page{11}
\note{Pagination de l'auteur : 8 (repassé).}

\struck{Proposition}

\struck{\emph{Lemme} a) Soit $a \in \mathcal{L}$ non \ill{} non minimal, et notons que $I_{<a}$ est défini. Alors $\operatorname{cosupp}^{\circ}(I_{<a}) = \{a\} \cup \operatorname{Omb}^{\circ}(I_{>a})$}
\note{Ces lignes sont barrées de trois longs traits obliques.}

\emph{Lemme.} $\mathfrak{M}$ est fermé dans $\widehat{\mathfrak{M}}$, i.e.\ pour $X, Y \in \widehat{\mathfrak{M}}$, pour qu'on ait $X \mathrel{|o|} Y$, il faut que pour tout $X' \mathrel{\mathring{\ll}} X$, $Y' \mathrel{\mathring{\ll}} Y$, $X', Y' \in \mathfrak{M}$, on ait $X' \mathrel{|o|} Y'$.

\marginal{NB C'est là le \uncertain{reflet}, \uncertain{dans} la \uncertain{situation}, de \ill{} $Y' \mathrel{\mathring{\ll}} X$ \uncertain{pour} $Y$ \ill{} \note{écrit en oblique dans la marge gauche, lecture très incertaine}}

\struck{Supp} Il n'y a de question que \struck{\ill{} si $X$ est un} \struck{\uncertain{\ill{} des deux éléments} $X$, $Y$ est} \uncertain{un} $I_{<a}$ \uncertain{ou un} $I_{>b}$. Supposons p.~ex.
\[
  X = I_{<a} ,
\]
\struck{et supposons d'abord} $Y \in \mathfrak{M}$. Alors

\struck{\ill{}} $X \mathrel{|o|} Y$ ssi $I_{<a} < Y$ i.e.\ $a \leq Y$, \struck{\ill{} \uncertain{montrons} \ill{} qu'il \uncertain{suffit} $X' \ll I_{<a}$,} \struck{\uncertain{ssi} $\forall\, X' \mathrel{\mathring{\ll}} I_{<a}$, \uncertain{on ait} $X' \mathrel{|o|} Y$.} \struck{$X' \leq a$, \uncertain{alors}}

Prenons d'abord $X' = I_{a',a}$, \uncertain{avec} $a' < a$ (\uncertain{tel} \uncertain{existe}, car $a$ non minimal). On aura $Y \mathrel{|o|} I_{a',a}$ donc $Y \trianglerighteq a$ \emph{ou} $Y \leq a'$. Mais si on avait $Y \leq a'$ ($< a$), on aurait $Y \mathrel{\mathring{\ll}} X = I_{<a}$ et prenant $X' = Y$, on trouve $Y \mathrel{|o|} Y$, absurde.
\note{Dans $X' = I_{a',a}$ l'indice est récrit sous un premier indice biffé. Le signe lu $\trianglerighteq$ dans « $Y \trianglerighteq a$ » est incertain (peut-être $\geq$).}

\struck{Si on prend \ill{}, pour $X'$ fixé, que}

\struck{$X' \mathrel{|o|} Y$ ssi} \struck{\uncertain{condition} avec \ill{}),}

\emph{Cor.} Les \uncertain{éléments} \struck{\uncertain{condition}} \uncertain{de} $\Sigma_{\widehat{\mathfrak{M}}}$ \uncertain{sont} \uncertain{canoniquement} \ill{} à l'un des supports $\Sigma_{\mathfrak{M}}$.
\note{Le corollaire est écrit au bas de la feuille, serré ; lecture incertaine.}

\page{12}
\note{Pagination de l'auteur : 9.}

Mais \uncertain{donc} à l'un des supports \emph{pour} $\mathfrak{M}$ qui \uncertain{\ill{}} \uncertain{s'intéresse} en fait !

Nous allons : \uncertain{présentement} travailler avec des figures et préfigures finies de $\widehat{\mathfrak{M}}$ (ou $\widehat{\mathfrak{M}}$-préfigures, ou préfigures \uncertain{généralisées}), i.e.\ les parties finies totalement ordonnées (et \uncertain{\ill{}}) de $\widehat{\mathfrak{M}}, \leq$. Pour une telle partie $\mathfrak{F}$, on introduit
\[
  T = \delta \mathfrak{F} \in \operatorname{Drap}^{*}(\mathcal{L}) ,
\]
\add{\ill{}} \uncertain{non vide} si $\mathfrak{F}$ est \uncertain{non vide}. \struck{Ainsi} On introduit les \uncertain{\emph{structures ensemblistes}} (\uncertain{généralisées}) d'un drapeau
\[
  T = \{ t_1 < t_2 < \cdots < t_n \}
\]
de $\mathcal{L}$, ainsi : \struck{\ill{}}
\[
  [I_{<t_1}],\ t_1,\ I_{t_1,t_2},\ t_2,\ I_{t_2,t_3},\ \ldots,\ t_n,\ [I_{>t_n}]
\]
\note{À côté de la marge gauche, en oblique : « \ill{} si \ill{} $b$ \ill{} \uncertain{non vide} » ; lecture incertaine.}
où les termes $I_{<t_1}$ (\uncertain{resp.}\ $I_{>t_n}$) ne figurent que si (\uncertain{si} $t_1$ non minimal, il n'existe \ill{}) \uncertain{et} \uncertain{seulement} \ill{} $t_1$ non minimal, et ne figurent que si $t_n$ non maximal (\uncertain{idem}\ldots).

Ce sont les éléments de la plus grande $\widehat{\mathfrak{M}}$-préfigure \add{possible} $\widehat{\mathfrak{F}}$, telle que $\delta\widehat{\mathfrak{F}} = T$. \struck{$\widehat{\mathfrak{M}}$-préfigure} \uncertain{Le} \uncertain{\ill{}} \uncertain{de} \add{$\mathfrak{F} \neq \emptyset$} \uncertain{revient} à celle de $\mathfrak{F}$, \uncertain{des} drapeaux $T$, plus un ensemble fini de \uncertain{structures} \uncertain{ensemblistes} (généralisées) \uncertain{\ill{}} de \uncertain{propriétés} \uncertain{suivantes} :

\note{La dernière ligne de la page, de lecture incertaine : « $\forall\, t \in T$, \ill{} $t \in \mathfrak{F}$, \emph{ou} $t \lhd X$ \uncertain{avec} $X \in \mathfrak{F}$ ».}
\[
  \forall\, t \in T, \quad t \in \mathfrak{F}, \ \text{ou} \ t \lhd X \ \text{avec} \ X \in \mathfrak{F} .
\]

\page{13}
\note{Pagination de l'auteur : 10.}

Ceci posé, on définit la \struck{préfigure} \emph{complémentaire} $\mathfrak{F}^{*}$ de $\mathfrak{F}$ comme
\[
  \mathfrak{F}^{*} = \widehat{\mathfrak{F}} \smallsetminus \mathfrak{F}
\]
$=$ ens.\ des \struck{\uncertain{intervalles}} \add{\uncertain{structures}} \uncertain{ensemblistes} généralisées de $\delta\mathfrak{F}$, qui ne sont pas $\in \mathfrak{F}$.

On aura
\[
  \delta \mathfrak{F}^{*} = \delta \mathfrak{F} \smallsetminus \mathfrak{F}_{\mathrm{red}} , \quad \text{soit } T'
\]
où $\mathfrak{F}_{\mathrm{red}}$ est l'ens.\ des sommets \uncertain{redondants} de $\mathfrak{F}$, i.e.
\[
  \mathfrak{F}_{\mathrm{red}} = \{ s \in \mathfrak{F}_0 \mid \operatorname{ord}(s, \mathfrak{F}) = 2 \} .
\]
\note{L'indice de $\mathfrak{F}_{\mathrm{red}}$ est écrit d'abord « kered » (?) puis « red » ; lecture incertaine.}

\struck{On \ill{}} Il y a lieu ici de faire la petite \uncertain{digression} sur la notion des sommets \uncertain{dits} d'une $\widehat{\mathfrak{M}}$-préfigure (\uncertain{cinq} types : sommets-bord propres et \uncertain{impropres}, sommets isolés, sommets intérieurs propres et impropres (\uncertain{ceux-là} \uncertain{seront} \uncertain{les} redondants) --- cf.\ p.~110, 111. On trouve \struck{\uncertain{aussi}} \add{\uncertain{ainsi}} que $\mathfrak{F}^{*}$ n'a pas de sommets redondants, \uncertain{donc} $\delta\mathfrak{F}^{**} = \delta\mathfrak{F}^{*}$.
\note{Les « p.~110, 111 » sont des pages de l'auteur, extérieures à ce lot.}

On obtient
$\mathfrak{F}^{**} = $ ens.\ des structures ensemblistes $X$ de $T' = \delta\mathfrak{F} \smallsetminus \mathfrak{F}_{\mathrm{red}}$ telles qu'il existe \struck{\ill{}} $Y \in \mathfrak{F}$ avec $Y \ll X$.
On dit que $\mathfrak{F}^{**}$ s'obtient à partir de $\mathfrak{F}$

\page{14}
\note{Pagination de l'auteur : 11.}

en « effaçant les sommets redondants », obtenus par « composition des \struck{\uncertain{segments}} intervalles adjacents aux sommets redondants ». On a
\[
  (\mathfrak{F}^{**})_0 = \mathfrak{F}_0 \smallsetminus \mathfrak{F}_{\mathrm{red}}
\]
$(\mathfrak{F}^{**})_1 = $ \struck{composés} ens.\ des intervalles \struck{\uncertain{obtenus}} \add{\uncertain{obtenus}} par composition des \struck{\uncertain{suites}} drapeaux maximaux \emph{connexes} contenus dans $\mathfrak{F}$,
i.e.\ de la forme
\[
  X_1,\ x_1,\ X_2,\ x_2,\ \ldots,\ x_{k-1},\ X_k
\]
avec \struck{$x_1 \lhd X_1, X_2,$}
\[
  X_1 < x_1 < X_2 < x_2 < \cdots < x_{k-1} < X_k
\]
\[
  \text{et} \quad x_1 \lhd X_1, X_2 ,\quad x_2 \lhd X_2, X_3 ,\ \ldots,\ x_{k-1} \lhd X_{k-1}, X_k .
\]
Elles sont en corr.\ 1-1 avec les composantes connexes de $\mathfrak{F}$, lesquelles sont définies par \uncertain{l'adjonction} \uncertain{éventuelle} d'un $x_0 \lhd X_1$, ou d'un $x_k \lhd X_k$, ou des deux.
\note{La deuxième égalité est écrite en trois lignes à droite du signe $=$ ; « composés » est biffé, un mot au-dessus de « obtenus » est biffé.}

\struck{\ill{}} Ceci posé, \uncertain{étant donné} une préfigure $\mathfrak{F}$, \uncertain{notons}
\[
  \operatorname{Omb}^{\circ}_{\mathfrak{M}}(\mathfrak{F}) = \bigcup_{X \in \mathfrak{F}} \operatorname{Omb}^{\circ}_{\mathfrak{M}}(X)
  = \{ Z \in \mathfrak{M} \mid \exists\, X \in \mathfrak{F},\ Z \mathrel{\mathring{\ll}} X \} .
\]
\note{Devant la réunion, un signe biffé.}

Ceci posé, le résultat clef est le suivant :
\note{L'énoncé suit sur la feuille suivante (p.~15).}

\page{15}
\note{Pagination de l'auteur : 12.}

\emph{Th.} Soit $\mathfrak{F}$ une $\widehat{\mathfrak{M}}$-préfigure \add{non vide}, \struck{et supposons \ill{} a) $\operatorname{or}\delta\mathfrak{F}$} \struck{\ill{} soit \ill{} a) $\operatorname{or}\delta\mathfrak{F}$}

\struck{Soit $a = \operatorname{or}\delta\mathfrak{F}$, $b = \operatorname{ex}\delta\mathfrak{F}$, et supposons que} \struck{a) si $a$ est minimal, i.e.\ un plus petit élément, et \add{si $\mathcal{L}_{<a} \in \mathfrak{F}$ et}} \struck{b) si $b$ est maximal, i.e.\ un plus grand élément \add{si $\mathcal{L}_{>b} \in \mathfrak{F}$ et}} \struck{(Ce qui \uncertain{résulte} \add{\uncertain{aussi}} de l'hypothèse $\mathrm{H}_{\mathcal{L}}$ : \ill{} \struck{\ill{}} $\mathcal{L}$ a un plus petit élément, \uncertain{ou bien} pour tout $a \in \mathcal{L}$, $\mathcal{L}_{<a} \neq \emptyset$, et $\mathcal{L}$ a un plus grand élément, \uncertain{ou} pour tout $a \in \mathcal{L}$, $\mathcal{L}_{>a} \neq \emptyset$.)}
\note{Toutes ces hypothèses sont barrées de deux longs traits obliques ; elles ont été abandonnées, et le théorème reste énoncé pour toute $\widehat{\mathfrak{M}}$-préfigure non vide.}

Alors
\[
  \operatorname{cosupp}^{\circ}_{\mathfrak{M}}(\mathfrak{F})
  = \operatorname{cosupp}^{\circ}\bigl(\operatorname{Omb}^{\circ}_{\mathfrak{M}}(\mathfrak{F})\bigr)
  = \operatorname{Omb}^{\circ}_{\mathfrak{M}}(\mathfrak{F}^{*}) .
\]

\emph{Corollaire 1.} Sous la condition $\mathrm{H}_{\mathcal{L}}$, \uncertain{plus} \uncertain{simplifiée} :
\[
  \operatorname{supp}^{\circ}_{\mathfrak{M}}(\mathfrak{F})
  = \operatorname{supp}^{\circ}_{\mathfrak{M}}\bigl(\operatorname{Omb}^{\circ}_{\mathfrak{M}}(\mathfrak{F})\bigr)
  = \operatorname{Omb}^{\circ}_{\mathfrak{M}}(\mathfrak{F}^{**}) .
\]
\note{Lecture incertaine des mots qui suivent « $\mathrm{H}_{\mathcal{L}}$ » ; les deux égalités sont celles de la page.}

\emph{Corollaire 2.} \struck{Soit (Sous la condition)} Soit $T \in \operatorname{Drap}(\mathcal{L})$. Considérons l'ens.\ des figures telles que
\note{La page s'arrête là ; le reste de la feuille est blanc.}

\page{16}
\note{Pagination de l'auteur : 13. La phrase du corollaire 2 (p.~15) n'est pas reprise ; la page s'ouvre sur un mot isolé.}

\uncertain{C.~facile}.

\emph{Remarques}, sur l'introduction des $I_{<a}$, $I_{>b}$.

\emph{Autre présentation.}

$1^{\circ}$) Pour y comprendre quelque chose, commençons avec le cas d'un « \struck{\ill{}} \emph{segment} \emph{ordonné} » ou « \emph{segment orienté combinatoire} » (que j'appelais aussi « \emph{pseudo-segment} », car on s'en assure par la suite en vérifiant les hypothèses), c'est la \emph{système}
\[
  \{ I, \omega_0, \omega_1 \} \ \text{ou} \ (I, \partial I) , \quad \partial I = \{\omega_0, \omega_1\}
\]
\add{(\uncertain{correspondant} au segment orienté)} d'un ensemble ordonné \struck{$I$} \struck{\ill{}} d'un plus petit élément $\omega_0$ ou \emph{origine}, et d'un plus grand élément $\omega_1$ ou \emph{extrémité}, avec
\[
  \omega_0 < \omega_1 \quad \text{i.e.} \quad \operatorname{card} I \geq 2 .
\]
On \uncertain{exclut} par là le cas où $\operatorname{card} I = 1$, i.e.
\[
  |I|^{\circ} = \emptyset , \quad \text{où} \quad |I|^{\circ} = I \smallsetminus \delta I \subset I .
\]
\note{Le premier « $|I|^{\circ}$ » est récrit sur un premier essai.}

Le cas est celui des \emph{segments} \emph{indivisibles}, qui donnent lieu à une \uncertain{magasin} qui est \add{\uncertain{souvent}} \struck{\ill{}} \emph{\uncertain{singulier}} \uncertain{simpliciale}, celle des simplexes orientés-type de dimension $1$, noté $\Delta^{1}$.
\marginal{C'est en vertu de la \uncertain{singularité} \ill{} \note{écrit en oblique dans la marge gauche ; lecture incertaine}}

Dans ce cas, on introduit pas des structures de la forme $I_{<a}$, $I_{>b}$. La notion de préfigure complémentaire $\mathfrak{F}^{*}$ d'une préfigure $\mathfrak{F}$ se développe, en restant dans le cadre des \uncertain{\ill{}} $a$ ($a \in I$) et $I_{a,b}$ ($a < b$, $a, b \in I$). Les développements \uncertain{concernant} la \uncertain{dite} \uncertain{\ill{}},

\page{17}
\note{Pagination de l'auteur : 14.}

des préfigures et des figures, leurs $\delta\mathfrak{F}$, $\partial\mathfrak{F}$, \uncertain{restent} \uncertain{valables} \ill{}, à un détail près, concernant \add{définitions et terminologie concernant les} les \emph{sommets} \emph{extérieurs} \struck{de $\mathfrak{F}$} ou « \uncertain{internes} » de $\mathfrak{F}$, propres ou impropres. En plus des sommets internes propres précédemment \uncertain{définis} (les $s \in \delta\mathfrak{F}$ t.q.\ $\operatorname{ord}(s, \mathfrak{F}) = 2$), \add{i.e.\ $s \in \partial\mathfrak{F}$ \ill{}} il faut \uncertain{joindre} les \struck{$s \in \delta\mathfrak{F}$} $s \in \partial I$ tels que $\operatorname{ord}(s, \mathfrak{F}) = 1$, (\emph{NB} si $s \notin \partial I$, on a toujours $\operatorname{ord}(s, \mathfrak{F}) \in \{0, 1\}$, à l'exclusion du \uncertain{cas} \uncertain{de} \uncertain{l'unique} \uncertain{élément} \uncertain{\ill{}} \uncertain{dans} \uncertain{\ill{}} $\mathfrak{F}$ ou \uncertain{bien} $\{X\}$, où $X \in \mathfrak{M}_1$, est \uncertain{contractile} \uncertain{\ill{}} \uncertain{\ill{}} $\mathfrak{F}$ de la forme \ill{} \uncertain{\ill{}} les \uncertain{\ill{}} de $\delta I$). On dit qu'un tel sommet est \struck{\ill{}} \emph{redondant} s'il est propre, \uncertain{\ill{}} \uncertain{impropre} si \uncertain{\ill{}} pour les sommets redondants proprement dits) il n'est pas dans $\delta\mathfrak{F}^{*}$ ; dans le cas inverse, il est dit sommet \uncertain{lacunaire}. Ceci dit, \ill{}.
\note{Toute la moitié supérieure de cette page est écrite vite ; la lecture de la parenthèse « NB \ldots » est très incertaine.}
\[
  \mathfrak{F}^{*}_{\mathrm{is}} = \mathfrak{F}_{\mathrm{lac}} , \qquad \mathfrak{F}^{*}_{\mathrm{lac}} = \mathfrak{F}_{\mathrm{is}}
\]
\[
  \mathfrak{F}_{\mathrm{red}} = \delta\mathfrak{F} - \delta\mathfrak{F}^{*}
\]
\[
  \delta\mathfrak{F} = \underbrace{\mathfrak{F}_{\mathrm{red}} \cup \mathfrak{F}_{\mathrm{lac}}}_{\mathfrak{F}_{\mathrm{int}}}
  \cup \mathfrak{F}_{\mathrm{is}} \cup \partial\mathfrak{F}
\]
\[
  \partial\mathfrak{F} = \underbrace{\partial_0\mathfrak{F}}_{\mathfrak{F}_0 \cap \partial\mathfrak{F}} \amalg \partial'_0\mathfrak{F}
\]
\note{Sous la dernière décomposition, de sa main : « $\partial_0\mathfrak{F}$ : sommets-bord propres », « $\partial'_0\mathfrak{F}$ : somm.\ b.\ impropres ».}

\marginal{Écrire \uncertain{sommets} \emph{disjoints} : \ill{}
\[
  \delta\mathfrak{F} = \delta\mathfrak{F} = \mathfrak{F}_{\mathrm{red}} \cup \mathfrak{F}_{\mathrm{is}} \cup \mathfrak{F}_{\mathrm{lac}} \cup \{\partial^{*}\mathfrak{F} \ldots
\]
où $\partial^{*}\mathfrak{F} = \partial\mathfrak{F} \smallsetminus \partial\mathfrak{F} \cap \partial I$ (sommets-bords non extrémaux)}
\note{Ajouté en bas à gauche, en partie récrit ; dans la première formule, des indices sont noircis.}

\uncertain{Les} \uncertain{\ill{}} disjoints, \uncertain{sauf} pour les \uncertain{deux} \add{propres} derniers termes, \uncertain{\ill{}} $\partial\mathfrak{F}$ et les deux premiers, formant $\mathfrak{F}_{\mathrm{int}}$, car
\[
  \partial\mathfrak{F} \cap \mathfrak{F}_{\mathrm{int}} = \partial I \cap \text{\struck{$\partial I$}} (\mathfrak{F}_0 \smallsetminus \mathfrak{F}_{\mathrm{is}})
\]
les points de $\delta\mathfrak{F}$ qui sont internes sans \uncertain{être} \uncertain{dans} $\partial I$, sans être des points isolés \struck{\uncertain{propres}} qui sont dans $\partial I$ de $\mathfrak{F}$.

\page{18}
\note{Pagination de l'auteur : 15.}

Du coup, l'intérêt de l'introduction de $\partial\mathfrak{F}$ \uncertain{devient} (les sommets-bord de $\mathfrak{F}$) devient plus problématique, d'autant plus que \struck{\uncertain{ces} sommets} \add{l'ensemble} \struck{\uncertain{bords}} \struck{de $\mathfrak{F}$} \uncertain{\ill{}} \uncertain{\ill{}} $\delta\mathfrak{F} = \partial\mathfrak{F}^{*}$. \struck{\uncertain{À vrai}}

\struck{En fait, on a}
\[
  \text{\struck{$\partial\mathfrak{F}^{*} \cap \delta\mathfrak{F} = \partial\mathfrak{F} \smallsetminus \partial\mathfrak{F} \cap \partial I$}}
\]
\[
  \text{\struck{$\partial\mathfrak{F} = (\partial\mathfrak{F}^{*} \cap \partial\mathfrak{F}) \amalg \partial\mathfrak{F} \cap \partial I$}}
\]
\[
  \text{\struck{$\partial\mathfrak{F} = (\delta\mathfrak{F}^{*} \cap \partial\mathfrak{F}) \amalg \mathfrak{F}_{\mathrm{is}} \cap \partial I$}}
\]
\note{Ces trois formules sont barrées de longs traits obliques ; sous le second terme de la deuxième, « $= \mathfrak{F}_{\mathrm{int}} \cap \partial I$ ».}

On a
\[
  \left.
  \begin{aligned}
    \mathfrak{F}_{\mathrm{is}} &= \mathfrak{F}^{*}_{\mathrm{lac}} \\
    \mathfrak{F}_{\mathrm{lac}} &= \mathfrak{F}^{*}_{\mathrm{is}}
  \end{aligned}
  \right\}
\]
\uncertain{sommets} lacunaires et s.\ isolés de $\mathfrak{F}$ et de $\mathfrak{F}^{*}$ s'échangent.
\[
  \mathfrak{F}_{\mathrm{red}} = \delta\mathfrak{F} \smallsetminus \delta\mathfrak{F} \cap \delta\mathfrak{F}^{*}
  \qquad \text{NB}\ \mathfrak{F}^{*}_{\mathrm{red}} = \partial I \smallsetminus (\partial I \cap \partial\mathfrak{F}) \ldots
\]
\[
  = \underbrace{(\delta\mathfrak{F} \smallsetminus \delta\mathfrak{F}^{**})}_{\text{sommets évanescents}}
  \amalg \underbrace{(\partial I \cap \partial\mathfrak{F})}_{\text{sommets redondants}}
\]
\[
  \partial^{*}\mathfrak{F} = \partial^{*}\mathfrak{F}^{*} \qquad \text{où} \quad
  \partial^{*}\mathfrak{F} = \partial\mathfrak{F} \smallsetminus \partial\mathfrak{F} \cap \partial I
\]
\note{Le « $=$ » écrit au-dessus de « NB » relie peut-être la seconde ligne à la première ; les accolades sous la deuxième ligne portent les mots « sommets évanescents » et « sommets redondants ».}

Ainsi dans la décomposition canonique
\[
  \delta\mathfrak{F} = \mathfrak{F}_{\mathrm{red}} \cup \mathfrak{F}_{\mathrm{is}} \cup \mathfrak{F}_{\mathrm{lac}} \cup \partial^{*}\mathfrak{F}
\]
les $2^{\mathrm{e}}$ et $3^{\mathrm{e}}$ termes s'échangent par dualité, le troisième \uncertain{\ill{}} \uncertain{invariant} par dualité, le premier \uncertain{disparaît} par dualité --- mais \struck{\ill{}} la partie \struck{$\partial\mathfrak{F} \cap \partial I$} $\mathfrak{F}_{\mathrm{red}} \cap \partial I$ \struck{\ill{}} \uncertain{réapparaît} par bidualité : \struck{sommets red.} \uncertain{Donc} le terme « redondant » \uncertain{pour} ces sommets n'est pas tellement adéquat, \struck{\ill{}} \struck{On pourrait les appeler} \uncertain{si} on adopte la terminologie \uncertain{sommets} \emph{évanescents} $\mathfrak{F}_{\mathrm{év}}$ \uncertain{et} \uncertain{somm.\ red.}\ $\mathfrak{F}_{\mathrm{red}}$ \uncertain{\ill{}} on pourrait appeler ces sommets semi-évanescents. On trouve alors une décomposition
\note{Les deux dernières lignes sont de lecture incertaine.}

\page{19}
\note{Pagination de l'auteur : 16.}

disjointe \uncertain{canonique} de
\[
  \widehat{\delta}\mathfrak{F} \overset{\mathrm{def}}{=} \delta\mathfrak{F} \cup \partial I
  = \delta\mathfrak{F} \cup \underbrace{(\partial I \smallsetminus \partial I \cap \delta\mathfrak{F})}_{= \mathfrak{F}^{*}_{\mathrm{sév}}}
\]
$\mathrm{sév} = $ semi-évanescents, ainsi :
\[
  \widehat{\delta}\mathfrak{F} = \mathfrak{F}_{\mathrm{is}} \cup \mathfrak{F}_{\mathrm{lac}} \cup \mathfrak{F}_{\mathrm{év}} \cup \mathfrak{F}_{\mathrm{sév}} \cup \mathfrak{F}^{*}_{\mathrm{sév}} \cup \partial^{*}\mathfrak{F}
\]
\note{Sous la formule, de sa main : une double flèche courbe relie $\mathfrak{F}_{\mathrm{is}}$ et $\mathfrak{F}_{\mathrm{lac}}$, une autre $\mathfrak{F}_{\mathrm{sév}}$ et $\mathfrak{F}^{*}_{\mathrm{sév}}$ ; une flèche descendante sous $\mathfrak{F}_{\mathrm{év}}$ porte « disparaît par dualité » ; une flèche en boucle sur $\partial^{*}\mathfrak{F}$. Sous $\mathfrak{F}^{*}_{\mathrm{sév}}$ : « $= \partial I \smallsetminus \partial I \cap \delta\mathfrak{F}$ » (une première écriture biffée), et plus bas une formule biffée par des hachures, « \ill{} $\mathfrak{F}^{*} \cap \delta\mathfrak{F}$ ».}

où
\[
  \mathfrak{F}_{\mathrm{is}} = \mathfrak{F}^{*}_{\mathrm{lac}} , \quad
  \mathfrak{F}_{\mathrm{lac}} = \mathfrak{F}^{*}_{\mathrm{is}} , \quad
  \partial^{*}\mathfrak{F} = \partial^{*}\mathfrak{F}^{*} , \quad
  \mathfrak{F}^{*}_{\mathrm{év}} = \mathfrak{F}^{**}_{\mathrm{év}} = \emptyset .
\]
\note{Dans les deux dernières égalités, des indices sont récrits sur d'autres, noircis ; « év » est la lecture la plus probable.}

Pour les $s \in \partial I$, il y a quatre positions possibles par \uncertain{rapport} à $\mathfrak{F}$ :
\begin{itemize}
  \item a) $s \in \mathfrak{F}_{\mathrm{is}} \iff s \in \mathfrak{F}^{*}_{\mathrm{lac}}$
  \item b) $s \in \mathfrak{F}_{\mathrm{sév}} \iff s \notin \delta\mathfrak{F}^{*}$
  \item a$'$) $s \in \mathfrak{F}_{\mathrm{lac}} \iff s \in \mathfrak{F}^{*}_{\mathrm{is}}$
  \item b$'$) $s \notin \delta\mathfrak{F} \iff s \in \mathfrak{F}^{*}_{\mathrm{sév}}$
\end{itemize}
\[
  \mathfrak{F}_{\mathrm{sév}} = \{ s \in \partial I \mid s \notin \delta\mathfrak{F}^{*} \}
  = \partial I \smallsetminus \partial I \cap \delta\mathfrak{F}^{*}
\]
\[
  \mathfrak{F}^{*}_{\mathrm{sév}} = \{ s \in \partial I \mid s \notin \delta\mathfrak{F} \}
  = \partial I - \partial I \cap \delta\mathfrak{F}
\]

\marginal{Écrit en oblique dans la marge gauche : « Cor.~1 \struck{\ill{}} $= \operatorname{supp}^{\circ}(\operatorname{Omb}^{\circ}\mathfrak{F}) = \operatorname{Omb}^{\circ}\mathfrak{F}^{**}$ » ; « Cor.~2 $\operatorname{supp}^{\circ}\mathfrak{F} = $ \ill{} ».}

On trouve alors
\[
  \boxed{\operatorname{cosupp}^{\circ}(\mathfrak{F}) = \operatorname{cosupp}^{\circ}\bigl(\operatorname{Omb}^{\circ}(\mathfrak{F})\bigr) = \operatorname{Omb}^{\circ}(\mathfrak{F}^{*})}
\]

On pourrait appeler \emph{parties contractibles} de $\mathfrak{M}$ les parties de la forme $\operatorname{Omb}^{\circ}(\mathfrak{F})$, $\mathfrak{F}$ préfigure. \struck{Il y a} L'application
\[
  \mathfrak{F} \longmapsto \operatorname{Omb}^{\circ}(\mathfrak{F}) \qquad \operatorname{Préfig}(\mathfrak{M}) \to \operatorname{Cons}(\mathfrak{M})
\]
est bijective. En effet, on récupère $\mathfrak{F}$
\note{Le nom « $\operatorname{Cons}$ » de l'ensemble but est de lecture incertaine.}

\page{20}
\note{Pagination de l'auteur : 17.}

à partir de $\operatorname{Omb}^{\circ}(\mathfrak{F}) = \Sigma$ par
\begin{itemize}
  \item $\mathfrak{F}_1 = $ ensemble des éléments de \struck{$\operatorname{Omb}^{\circ}(\mathfrak{F})$} $\Sigma$ \uncertain{maximaux} pour $\mathring{\ll}$, qui sont dans $\mathfrak{M}_1$ ;
  \item \struck{$\mathfrak{F}_0 = \{ s \in \ldots \mid$}
  \item $\mathfrak{F}_{\mathrm{is}} = $ ensemble des éléments de $\Sigma$, maximaux pour $\mathring{\ll}$, qui sont dans $\mathfrak{M}_0$ ;
  \item $\mathfrak{F}_0 \smallsetminus \mathfrak{F}_{\mathrm{is}} = \Bigl( \bigcup_{X \in \mathfrak{F}_1} \partial X \Bigr) \cap \Sigma$
    $= \{ s \in \Sigma \cap \mathfrak{M}_0 \mid \exists\, X \in \mathfrak{F}_1 \ \text{t.q.}\ s \lhd X \}$.
\end{itemize}
\note{Ces quatre lignes sont réunies sur la page par une grande accolade.}
\note{Après le premier signe $=$ de la dernière ligne, une formule est biffée par des hachures (« $\{ s \in \Sigma \cap \mathfrak{M}_0 \}$ » ?).}

On trouve que le cosupport d'une partie contractible de $\mathfrak{M}$ est encore contractible.

\marginal{À gauche, en regard : « $\mathfrak{F}^{**}$ ».}

\struck{\ill{}} Et le passage au cosupport, \add{\uncertain{associé}} \add{\uncertain{du côté des}} préfigures, au passage à \struck{\uncertain{supposer}} la préfigure complémentaire. Le passage \uncertain{correspondant} \add{\uncertain{au passage}} à l'\emph{enveloppe} \emph{\uncertain{irrédondante}} \uncertain{de} la figure considérée, \struck{\ill{}} \add{\uncertain{associée}} $T \in \operatorname{Drap}(I)$, $T \supset \partial I$,
\[
  t_0 = \omega_0 < t_1 < \cdots < t_n = \omega_1
\]
(un drapeau), les structures \uncertain{ensemblistes} sont
\[
  \underset{\substack{\| \\ \omega_0}}{\{t_0\}},\ I_{t_0 t_1},\ t_1,\ I_{t_1, t_2},\ \ldots,\ I_{t_{n-1}, t_n},\ \underset{\substack{\| \\ \omega_n}}{t_n} .
\]
\note{Sous $t_n$, l'indice lu $\omega_n$ est probablement un $\omega_1$ ; on garde la page.}

\uncertain{Désignons} \uncertain{par} \emph{préfigure} \emph{subordonnée} à $T$, toute partie $\mathfrak{F}$ de cet ensemble \uncertain{des} drapeaux $T$, \uncertain{notée} $\mathfrak{F}_T$. On trouve dans \uncertain{ce} principe de \uncertain{dualité} \uncertain{précédente} \uncertain{ceci}, $\forall\, \mathfrak{F} \subset \mathfrak{F}_T$, on désigne par $\complement\mathfrak{F}$ son complémentaire dans $\mathfrak{F}_T$
\note{La fin de la page est d'une écriture rapide, la lecture des mots est très incertaine ; la phrase continue au-delà de ce lot.}

\end{document}
